% 第八章 展望与结论
\chapter{展望与结论}
\label{ch:outlook}

\section{研究总结}

本研究围绕有机光电存算一体（Compute-in-Memory, CIM）系统的算法—器件协同设计，构建了一套面向混合模拟/数字推理的完整仿真栈，系统评估了视觉Transformer（ViT）与大语言模型（LLM）KV缓存在有机模拟交叉阵列上的可行性、鲁棒性与能效边界。研究的核心贡献可概括为以下四个方面。

第一，建立了从器件物理到系统精度的一阶行为级仿真框架（第~\ref{chap:methodology}~章），涵盖差分对电导映射、多电平量化、D2D/C2C噪声、双指数保持衰减、光电前端补偿及HAT等模块。针对Tiny-ViT-5M，本研究将87.7\%参数映射至模拟域，覆盖注意力投影与前馈等密集算子；动态注意力运算保留数字域执行，构成系统能效的"模拟天花板"。

第二，表明Transformer对模拟非理想性的敏感度高於CNN，尺度恢复（scale recovery）对维持激活动态稳定不可或缺（第~\ref{chap:methodology}~章）。

研究进一步发现"硬件实例过拟合"（第~\ref{ch:failure_modes}~章）：固定D2D掩膜使模型记忆特定制造偏差，面对全新实例时崩溃至~$10.00\%$。Ensemble HAT通过在轮次边界重采样D2D失配图，将器件失配转化为随机正则化项，将fresh-instance精度恢复至~$86.16\pm0.19\%$。

第三，通过多维物理应力测试量化了系统的生存边界（第~\ref{ch:failure_modes}~章与第~\ref{chap:deployment-envelope}~章）：Ensemble HAT在标准噪声体制下维持~$86.16\%$ fresh-instance精度，但ADC低于6-bit或比例噪声下出现刚性崩塌；CIFAR-100等细粒度任务获益最大，Flowers-102等小数据场景受floor效应限制。

第四，将研究视野扩展至LLM KV缓存场景（第~\ref{ch:work2_scope}~章）。Remote~107阶段性审计结果提示：全层模拟KV缓存在8-bit零噪声下已超出10\%困惑度退化门限（17.48 vs 15.68 PPL），而选择性终端层（last1）在fresh-D2D验证中表现稳定（PPL=18.42--18.60，变异系数$<0.3\%$），可作为有机CIM KV缓存的候选路径；该方向仍需manifest补全与最小重跑后锁定。

\section{未来研究方向}

未来工作首先需要把仿真框架推广至ImageNet-1K等大规模数据集，并探索DeiT-S、Swin-T等更大视觉Transformer的映射可行性。模型规模的扩展将考验阵列容量与片上带宽的三角平衡，数十亿参数模型的全映射将不可避免地走向多芯片铺排。

此外，LLM KV缓存场景下的选择性终端层HAT已在TinyLlama-1.1B上显示出初步可行信号，其向更大规模模型（如LLaMA-2-7B GQA）的扩展将验证有机CIM在分组查询注意力架构下的可行性，并为长上下文（32k+）场景中的保持时间假设提供关键数据。

更关键的限制来自物理实验闭环。本研究目前基于文献先验与行为级参数化模型进行仿真，尚未与真实有机光电器件的晶圆级测量数据闭环。未来的首要任务是建立从实验室表征到仿真参数的自动化转换管线：将多级编程I-V曲线的电导窗口与积分非线性、重复读写统计的C2C方差、跨芯片阵列的D2D分布、保持衰减痕迹的双指数或stretched-exponential拟合，以及光响应非线性测量结果，通过统一接口自动拟合为JSON器件画像并直接驱动训练与推理脚本。

具体而言，当前默认的双指数保持模型（$\tau_1=140$~ms, $\tau_2=610$~ms, $A_0=0.6$）需要替换为实际有机半导体材料（如DNTT、C8-BTBT及其衍生物）在不同温度与偏置条件下的测量值；电导积分非线性（INL）应通过实测电平码表而非均匀量化理想化来表征；而NL\_LTP/NL\_LTD的梯度缩放近似也应升级为脉冲级写verify动力学模型。只有在实测数据闭环后，仿真框架才能从行为级设计工具进化为真正的工艺协同设计（DTCO）平台。

系统级能效模型也需要从一阶估算走向外围电路显式建模。现有能耗模型将阵列到数字互连能耗笼统地归入SRAM读写项，未单独核算专用布线、数据编排（data marshaling）、模拟—数字接口以及层间尺度恢复乘法器的开销。未来需引入更精细的外围电路模型，包括：按阵列分块（tiling）计费的ADC/DAC转换次数、可编程跨阻放大器（TIA）增益校准的能耗与面积代价、以及不同精度ADC在各模拟子算子上的分层自适应配置。

此外，当前模型尚未包含温度效应对载流子迁移率、阈值电压漂移及有机半导体能带重整化的影响，而这些在有机器件中尤为显著且与保持衰减强耦合。建立温度敏感的系统级能耗—精度联合模型，并引入位置相关的IR压降与潜行电流（sneak-path）SPICE校准模型，是向实际芯片部署迈进的必要步骤。

在可靠性层面，当前映射假设一次性编程（one-time programming）后权重静态驻留，忽略了长期推理过程中的电导漂移、读取干扰与热老化。未来应研究闭环写验证（closed-loop write-verify）策略，即在编程阶段通过逐脉冲反馈将电导精确收敛到目标窗口，并在运行期间引入周期性刷新、自适应偏置补偿或动态余量缩放（dynamic margin scaling），以对抗状态相关的保持衰减。

对于KV缓存场景，光学刷新机制——利用光子调制恢复电导态而不消耗耐久预算——是有机OEC-RAM的独特优势。系统评估光学刷新的能效边界、局部可控性与刷新策略，将为分钟级会话长度的KV缓存管理提供定量依据。对于高可靠性边缘应用，还需探索模拟纠错码（analog ECC）、冗余阵列或三模冗余（TMR）方案。

最后，当单层阵列无法满足大模型参数量时，多芯片阵列铺排（tiling）与片间路由成为必然选择。未来的研究需解决模拟权重在多个物理芯片间的最优划分、片间模拟信号完整性（串扰、衰减与阻抗匹配）、以及数字控制逻辑的分布式调度与同步等问题。有机光电CIM的独特优势在于光调控的多模态编程能力，这为利用光互连实现三维堆叠或晶圆级键合（wafer-to-wafer bonding）提供了天然接口：光信号既可作为权重编程手段，也可作为片间高带宽、低电容的通信载体。

\section{结论}

有机光电存算一体技术为突破冯·诺依曼架构能效瓶颈提供了替代方案，但将现代深度学习推理映射至有机阵列并非简单的权重迁移，而需要算法、电路与器件的深度协同优化。忽视动态注意力数字开销、ADC量化阈值或硬件实例过拟合风险将导致实验室原型与真实部署之间的显著鸿沟。

研究表明，Tiny-ViT在标准器件波动下可保持$86.16\pm0.19\%$的fresh-instance精度，选择性终端层KV缓存（last1, PPL=18.42--18.60）则提供了一个仍待锁定的LLM推理候选方向。研究同时界定了物理边界：动态注意力的数字执行构成模拟天花板，系统能效提升被限制在约$11.45\times$；ADC分辨率存在6-bit刚性阈值；硬件实例过拟合与数据量floor效应限制了小数据场景适用性。材料层面的电导优化必须与读出精度、训练鲁棒性及数字—模拟分工协同考虑，方能实现有意义的端到端加速。本研究建立的仿真框架与开源工具链将为后续向真实器件闭环与芯片原型转化提供方法论支撑。
